\documentclass[11pt,a4paper]{article}
\usepackage[margin=2.5cm]{geometry}
\usepackage{amsmath,amssymb,amsthm,mathtools,booktabs,array,xcolor}
\usepackage[colorlinks=true,linkcolor=blue!60!black,citecolor=green!40!black,urlcolor=blue!60!black]{hyperref}
\newcommand{\La}{\textsf{\textbf{(a)}}}
\newcommand{\Lb}{\textsf{\textbf{(b)}}}
\newcommand{\Lc}{\textsf{\textbf{(c)}}}
\newcommand{\Ld}{\textsf{\textbf{(d)}}}
\newcommand{\script}[1]{\texttt{scripts/#1}}
\newtheorem{proposition}{Proposition}[section]
\theoremstyle{remark}\newtheorem{remark}[proposition]{Remark}
\newcommand{\cA}{\mathcal A}

\title{A candidate exact two-soliton S-matrix\\ for the bad Boussinesq equation}
\author{Project \emph{quantum\_boussinesq\_solitons} (Phase 2, next step 3)}
\date{3 November 2026}
\begin{document}\maketitle

\begin{abstract}
We construct a candidate exact quantum two-soliton S-matrix for $v_{tt}=v_{xx}+v_{xxxx}+3(v^2)_{xx}$ and test it against everything Phase~2 established. The candidate is a ratio of Gamma functions in the \emph{spectral-parameter difference} $\eta=\tfrac29(\ell_1-\ell_2)/\hbar$, times a power of a cross-ratio of the cubic's roots, times the exact topological factor:
\[
S_{\rm ov}=\frac{\Gamma(1+i\eta)}{\Gamma(1-i\eta)}\Big(\frac{e^2A_{12}}{\eta^2}\Big)^{i\eta}e^{\frac{i}{\hbar}\frac23M_1M_2(v_1-v_2)},\qquad
S_{\rm ho}=\frac{\Gamma(1-i\eta)}{\Gamma(1+i\eta)}\Big(\frac{A_{12}\eta^2}{e^2}\Big)^{i\eta}e^{\frac{i}{\hbar}\frac23M_1M_2(v_1-v_2)} .
\]
It reproduces the semiclassical phase $\delta$ (including the topological term) up to a constant Maslov-type phase $\pm\pi/2$; it is exactly unitary, real-analytic, and obeys the crossing relation exactly (this fixes $S_{\rm ho}$). On the physical sheet its only poles at breather kinematics are simple poles exactly at the Bogdan--Charkina Bohr--Sommerfeld points $I_{\rm br}=N\hbar$, $N\ge1$. Fused breather--soliton amplitudes are unitary with the right classical limit. Three other families are \emph{excluded} at leading order: Gamma functions of root differences, any product $F(\beta_1-\beta_2)G(\beta_1+\beta_2)$ (so the pure $\mathfrak{sl}_3$ ``mirror'' form, and quantum dilogarithms of $\beta_1\pm\beta_2$), and the pure exponential $e^{i\delta/\hbar}$ (no bound states). The candidate is a guess \Ld\ whose properties are derived \Lc; it is \emph{not} derived from a quantum theory. Its remaining ambiguities are all $O(\hbar^0)$ in the phase and are listed with the input that would fix them.
\end{abstract}
\tableofcontents

\section{Setting and requirements}
Conventions are those of \texttt{report.pdf}: $k=\sinh\beta$, $v=\epsilon\cosh\beta$, $\ell=\epsilon\cosh3\beta$, $M=2k$, $P=\tfrac23\epsilon k^3\cosh\beta$; particle~1 starts on the left. The semiclassical phase is \Lc
\begin{equation}\label{eq:delta}
\delta(1,2)=\tfrac29(\ell_1-\ell_2)\log A_{12}+T_{12},\qquad T_{12}=\tfrac23M_1M_2(v_1-v_2),
\end{equation}
with $S\simeq e^{i\delta/\hbar}$. An exact candidate must (i) reproduce~\eqref{eq:delta} as $\hbar\to0$, (ii) be unitary, (iii) obey the crossing-type relation $S_{\rm ov}(\beta_1,\beta_2+i\pi)S_{\rm ho}(\beta_1,\beta_2)=1$ (the exact version of $\delta_{\rm ov}(\beta_1,\beta_2+i\pi)=-\delta_{\rm ho}(\beta_1,\beta_2)$), (iv) have bound-state poles consistent with the string condition $\ell_1-\ell_2\in\mp\tfrac{9i}2\hbar\mathbb N$, equivalently with the Bogdan--Charkina (BC) spectrum $I_{\rm br}=N\hbar$ \La~\cite{BC}, and (v) give consistent fused amplitudes. Scattering is diagonal, so Yang--Baxter is automatic.

\paragraph{Roots of the cubic.} \Lc\ (\script{smatrix\_ansatz.py}) The three roots attached to a right-mover are
\[
\mu_j=\tfrac{i}{\sqrt3}\cosh\!\big(\beta+\tfrac{i\pi(2j+1)}3\big),\quad j=-1,0,1,
\]
roots of $\mu^3+\tfrac14\mu=\tfrac{i}{12\sqrt3}\ell$; the first two are the pair $\tfrac12(iv/\sqrt3\pm k)$ of Phase~2, whose cross-ratio is $A_{12}$. Consequently, for every root $\mu_{1a}$ of soliton~1,
\begin{equation}\label{eq:prod}
\prod_{b}(\mu_{1a}-\mu_{2b})=\tfrac{i}{12\sqrt3}(\ell_1-\ell_2).
\end{equation}
So $\ell_1-\ell_2$ is itself a product of root differences, and every function below is a product of powers of root differences.

\section{Families tested and excluded}\label{sec:excl}
Single-particle factors are a convention, so we compare the mixed derivative $\partial_{\beta_1}\partial_{\beta_2}$ of the leading phase with that of $\delta$, by least squares at 40 random overtaking points (\script{smatrix\_rejected\_forms.py}). \Lc
\begin{center}\small
\begin{tabular}{@{}p{6.3cm}p{5.4cm}p{2.6cm}@{}}\toprule
Family & Leading phase spans & Residual\\\midrule
$\prod\Gamma(1+ic\,x_{ab}/\hbar)^{n_{ab}}$, $x_{ab}=\mu_{1a}-\mu_{2b}$ (any complex $c,n$), $\times e^{iT/\hbar}$ & $\Re,\Im\ x\log x$, $x$, $T$ & $2\times10^{-2}$: excluded\\
$F(\beta_1-\beta_2)\,G(\beta_1+\beta_2)$: $\mathfrak{sl}_3$ Toda / sine-Gordon mirror blocks, quantum dilogarithms of $\beta_1\pm\beta_2$ & $f(D)+g(S)$, so $\partial_S\partial_D=0$ & $\partial_S\partial_D\delta\neq0$, e.g.\ $-0.55$ at $(S,D)=(1.4,0.6)$: excluded\\
$e^{i\delta/\hbar}$ exactly & $\delta$ & passes (i)--(iii), \emph{no poles}: excluded by (iv)\\
Control: $\Gamma$ in $\eta$ (Sec.~\ref{sec:cand}) & $(\ell_1-\ell_2)\log x_{ab}$, $T$ & $5\times10^{-16}$\\
\bottomrule\end{tabular}\end{center}
The obstruction to the first two is the same: $\delta$ has the \emph{uniform} weight $\ell_1-\ell_2$ in front of every logarithm, while Stirling asymptotics of $\Gamma(\ldots+iX/\hbar)$ give weight $X$ in front of $\log X$. The weight must therefore be the Gamma argument. This forces the Gamma functions to depend on $\eta\propto\ell_1-\ell_2$ only. \Lc
The quantum-dilogarithm case is excluded only for arguments that are functions of $\beta_1\pm\beta_2$ separately; a dilogarithm with mixed arguments was not tested.

\section{The candidate}\label{sec:cand}
Let $\eta=\tfrac29(\ell_1-\ell_2)/\hbar$ and $\cA_{12}=\dfrac{81A_{12}}{4(\ell_1-\ell_2)^2}$, so $A_{12}/\eta^2=\hbar^2\cA_{12}$. \Ld
\begin{equation}\label{eq:Sov}
S_{\epsilon\epsilon}(1,2)=\frac{\Gamma(1+i\eta)}{\Gamma(1-i\eta)}\,\big(e^2\hbar^2\cA_{12}\big)^{i\eta}\,e^{iT_{12}/\hbar}\qquad(\epsilon_1=\epsilon_2),
\end{equation}
\begin{equation}\label{eq:Sho}
S_{+-}(1,2)=\frac{\Gamma(1-i\eta)}{\Gamma(1+i\eta)}\,\Big(\frac{A_{12}\eta^2}{e^2}\Big)^{i\eta}\,e^{iT_{12}/\hbar}\qquad(\text{head-on}).
\end{equation}

\begin{proposition}[\Lc, \script{smatrix\_ansatz.py}]\label{prop:struct}
\begin{enumerate}
\item $\cA_{12}$ is finite and nonzero at threshold $\beta_1=\beta_2$, so~\eqref{eq:Sov} has no singularity there: the $\log(\ell_1-\ell_2)^2$ singularity of $\log A_{12}$ is carried by the Gamma ratio. Using~\eqref{eq:prod}, $\cA_{12}$ is a ratio of products of root differences.
\item $\log\cA_{12}=h(\beta_1+\beta_2)+g(\beta_1-\beta_2)$: the $\mathfrak{sl}_3$/sine-Gordon blocks of Phase~2 (Prop.~5.1 of \texttt{report.pdf}) survive, but only \emph{in the exponent}, multiplied by $\eta$.
\item $\Gamma(ix)/\Gamma(-ix)=-\Gamma(1+ix)/\Gamma(1-ix)$. Hence the Gamma part is that of the quantum Toda two-body amplitude $R=\Gamma(i\kappa)/\Gamma(-i\kappa)\hbar^{2i\kappa}$ (\texttt{toda\_route.pdf}), with $\kappa=\Delta p/\hbar\mapsto\eta$, and the Toda Lax eigenvalue $p\mapsto\tfrac29\ell$. That $p$ maps to $\ell$ in the phase weight was found in Phase~2b; that the Gamma structure survives the Boussinesq limit is \Ld.
\end{enumerate}
\end{proposition}

\section{Tests}
\subsection{Semiclassical limit, unitarity, crossing}
\begin{proposition}[\Lc, \script{smatrix\_semiclassical\_unitarity\_crossing.py}]\label{prop:tests}
For real rapidities, at three overtaking and three head-on points:
\begin{enumerate}
\item $\arg S-\delta/\hbar\to\tfrac{\pi}2\sigma\,\mathrm{sgn}\,\eta$ (mod $2\pi$), $\sigma=+1$ (same direction), $-1$ (head-on), with remainder $O(\hbar)$; at $\hbar=10^{-2},10^{-3},10^{-4}$ the remainder is e.g.\ $-6.4\times10^{-4},-6.4\times10^{-5},-6.4\times10^{-6}$. Without the factor $e^{iT/\hbar}$ the remainder is $O(1/\hbar)$, so the topological term is required.
\item $|S|=1$ and $S(1,2)S(2,1)=1$ to 30 digits.
\item $S_{\rm ov}(\beta_1,\beta_2+i\pi)\,S_{\rm ho}(\beta_1,\beta_2)=1$ to $10^{-27}$; at leading order this is $\delta_{\rm ov}(\beta_1,\beta_2+i\pi)+\delta_{\rm ho}=0$.
\item $S(\bar\beta_1,\bar\beta_2)=1/\overline{S(\beta_1,\beta_2)}$ at generic complex rapidities.
\end{enumerate}
\end{proposition}
Crossing is what fixes~\eqref{eq:Sho}: under $\beta_2\to\beta_2+i\pi$, $\eta_{\rm ov}\to\eta_{\rm ho}$ and $A_{\rm ov}\to1/A_{\rm ho}$ (Phase~2), so the inverse Gamma ratio and the power $(A\eta^2/e^2)^{i\eta}$ are forced. Using~\eqref{eq:Sov} for head-on scattering too would violate crossing. During the work we first wrote~\eqref{eq:Sho} with $(\eta^2/e^2A)^{i\eta}$; the semiclassical test failed, and crossing gave the correct form.

The constant $\pm\pi/2$ is a Maslov-type phase at $O(\hbar^0)$. The report listed a missing Maslov phase among its least-sure points; the candidate predicts one. \Ld

\subsection{Poles and the breather spectrum}
Take two right-movers at $\beta_{1,2}=a\mp ib$ with $\pi/3<b<\pi/2$ (right-moving breathers; left-moving ones are their parity images).

\begin{proposition}[\Lc, \script{smatrix\_poles\_bootstrap.py}]\label{prop:poles}
\begin{enumerate}
\item With this labelling $\operatorname{Im}(P_1-P_2)>0$: the relative momentum lies in the upper half plane, so this is the physical sheet for bound states of $S(1,2)$. Total $E$ and $P$ are real.
\item $i\eta=\tfrac49\sinh3a\sin3b/\hbar=-I_{\rm br}/\hbar$, with $I_{\rm br}=-\tfrac49\sinh3a\sin3b$ the BC adiabatic invariant. Hence $S(1,2)$ has simple poles (from $\Gamma(1+i\eta)$) \emph{exactly} at $I_{\rm br}=N\hbar$, $N=1,2,\ldots$: the BC Bohr--Sommerfeld spectrum, with their $N$ (not $N+\tfrac12$). Verified for $N\le4$ at $a=0.5,0.8$, $\hbar=0.1$ ($|S||\delta b|$ constant to $10^{-9}$).
\item There is no pole at $N=0$ (threshold $b=\pi/3$); there the Toda form $\Gamma(i\eta)$ would have one. This selects $\Gamma(1+i\eta)$.
\item The form that is regular at real threshold and the form with poles on the physical sheet coincide. This was not imposed and is a non-trivial consistency check.
\item At a pole $e^{i\eta\cdot2\pi i}=1$, so the residue modulus does not depend on the branch of $\log\cA_{12}$ (which lies on its cut, $\cA_{12}<0$, at breather kinematics). The residue phase does.
\end{enumerate}
\end{proposition}

Sample energies $E_N=E(\beta_1)+E(\beta_2)$ at $\hbar=0.1$, $a=0.8$: $0.8711, 0.9759, 1.0759, 1.1705$ for $N=1..4$ (\Lc, numerical; identical to BC's quantisation because the conditions coincide).

\subsection{Bootstrap: breather--soliton amplitudes}
\begin{proposition}[\Lc, \script{smatrix\_poles\_bootstrap.py}]\label{prop:boot}
Define $S_{B_N,3}=S(1,3)S(2,3)$ at the $N$-th pole. Then:
\begin{enumerate}
\item $|S_{B_N,3}|=1$ for real $\beta_3$ of either direction, and $\delta_{13}+\delta_{23}$ is real (the classical breather--soliton phase, by Hirota additivity).
\item $\arg S_{B_N,3}-(\delta_{13}+\delta_{23})/\hbar\to\pi$ as $\hbar=I_{\rm br}/N\to0$ at fixed $(a,b)$: two Maslov-type constants.
\item With $x_c=\tfrac i2(\eta_{13}+\eta_{23})$, the Gamma part telescopes to
\[
\frac{\Gamma(1+x_c-\frac N2)\Gamma(1+x_c+\frac N2)}{\Gamma(1-x_c+\frac N2)\Gamma(1-x_c-\frac N2)} .
\]
It has a simple pole at $x_c=-\tfrac N2$, double poles at $x_c=-\tfrac N2-1-m$, and the mirror-image zeros; $N-1$ pole--zero pairs cancel in between. This was checked numerically for $N=1,2,3$.
\end{enumerate}
\end{proposition}
\emph{Interpretation} \Ld: the simple pole is at $\ell_3=\ell_2$, $\ell_1-\ell_3=$ $N$-string, an exchange of soliton~3 with a constituent. The double poles extend the string, $\ell_3=\ell_2-\tfrac{9i}{2}(m+1)\hbar$, and would need a Coleman--Thun-type explanation. Whether these lie on the physical sheet of the three-body kinematics was \emph{not} determined, so the bootstrap is only partially tested. Breather--breather amplitudes follow by the same fusion but were not analysed.

\section{Summary of tests}
\begin{center}\small
\begin{tabular}{@{}p{4.3cm}p{8.0cm}p{2.3cm}@{}}\toprule
Requirement & Result for~\eqref{eq:Sov}--\eqref{eq:Sho} & Status\\\midrule
Semiclassical limit & $\delta/\hbar\pm\pi/2+O(\hbar)$, topological term needed & \Lc\ passed\\
Unitarity & exact & \Lc\ passed\\
Crossing & exact; fixes the head-on form & \Lc\ passed\\
Real analyticity & exact & \Lc\ passed\\
Poles vs string / BC & simple poles exactly at $I_{\rm br}=N\hbar$, physical sheet, none at $N=0$ & \Lc\ passed\\
Bootstrap & fused amplitudes unitary, right classical limit; pole interpretation open & \Lc\ partial\\
Derivation from a quantum theory & none & \Ld\\
\bottomrule\end{tabular}\end{center}

\section{Remaining ambiguities and what would fix them}
All tests constrain the phase only at orders $\hbar^{-1}$ (fixed by $\delta$) and at the pole positions. The following modifications pass every test above, or fail only the poles in a way the present data cannot decide. \Lc\ (that they pass) / \Ld\ (that this list is complete)
\begin{enumerate}
\item \textbf{CDD factors.} Any $f(1,2)$ with $f(1,2)f(2,1)=1$, $|f|=1$ for real rapidities, $f_{\rm ov}(\beta_1,\beta_2+i\pi)f_{\rm ho}(\beta_1,\beta_2)=1$, and $\log f=O(\hbar^0)$. Examples: $\prod_j\frac{\sinh\frac{\pi}{2}(\eta-i\alpha_j)}{\sinh\frac{\pi}{2}(\eta+i\alpha_j)}$, or $\frac{c+i\eta}{c-i\eta}$ with constant $c$. They change one-loop phases and add or remove poles; the latter could be tested against a quantum breather spectrum.
\item \textbf{Renormalisation of $\hbar$} in $\eta$, $\hbar\to\hbar_{\rm ren}=\hbar(1+c_1\hbar+\ldots)$, as in sine-Gordon $\beta^2\to\beta^2/(1-\beta^2/8\pi)$. This shifts the breather spectrum at $O(\hbar^2)$.
\item \textbf{Shift of the string}: $\Gamma(\tfrac12+i\eta)$ instead of $\Gamma(1+i\eta)$, i.e.\ $I_{\rm br}=(N-\tfrac12)\hbar$. This is excluded only if one trusts BC's integer quantisation, which is itself a leading-order Bohr--Sommerfeld statement.
\item \textbf{$O(\hbar)$ corrections to $\cA_{12}$} inside the power, and the branch of $\log\cA_{12}$ at complex rapidities (residue phases).
\item \textbf{The Maslov-type constant} $\pm\tfrac\pi2$, and whether the overall sign of $S$ is a convention.
\end{enumerate}

\paragraph{Input that would fix them.}
\begin{itemize}
\item \emph{One-loop soliton and breather spectrum} (DHN-type; next step 2 of the report). The $O(\hbar^0)$ correction to the breather energies $E_N$ fixes items 2--3. The one-loop two-soliton phase (from the phonon phase shifts on the reflectionless soliton background) fixes item~4 and tests the $\pm\pi/2$.
\item \emph{Small-amplitude Toda limit} (next step 1: the Fermi-edge double scaling $\beta\to0$, $\hbar\propto\beta^2$, of the quantum Toda Bethe ansatz). This gives an exact quantum S-matrix in the corner $\beta\to0$ against which~\eqref{eq:Sov}--\eqref{eq:Sho}, including possible CDD factors, can be compared directly. Note $\eta=O(1)$ there means $\hbar\propto\beta^4$ for overtaking ($\ell_1-\ell_2\approx\tfrac92(\beta_1^2-\beta_2^2)$) but $\hbar=O(1)$ for head-on ($\ell_1-\ell_2\approx2$). So the head-on Gamma ratio is never semiclassical in that corner. \Lc\ (elementary)
\item \emph{A quantum lattice regularisation} at the dispersion scale would be needed to derive the candidate rather than guess it; none is known (\texttt{toda\_route.pdf}).
\end{itemize}

\appendix\section{Scripts}
All are in \texttt{phase2/scripts}; each ends with \texttt{ALL CHECKS PASSED}.
\begin{itemize}
\item \texttt{smatrix\_ansatz.py}: definitions; roots, Eq.~\eqref{eq:prod}, Prop.~\ref{prop:struct}.
\item \texttt{smatrix\_semiclassical\_unitarity\_crossing.py}: Prop.~\ref{prop:tests}.
\item \texttt{smatrix\_poles\_bootstrap.py}: Props~\ref{prop:poles}--\ref{prop:boot}.
\item \texttt{smatrix\_rejected\_forms.py}: Section~\ref{sec:excl}.
\end{itemize}
\emph{Failures along the way:}
\begin{itemize}
\item The first root labelling mismatched the Phase~2 pair. It was fixed by matching to \texttt{consistency\_checks.py}.
\item The first head-on form had the wrong sign of the $A$-power; it failed the semiclassical and crossing tests and was corrected as described above.
\item A threshold-size tolerance was set too tightly ($|\cA|>10^{-3}$, whereas $\cA\approx2\times10^{-4}$ at $\beta=1.6$).
\item $|S|$ at complex breather kinematics is very large ($\sim10^{23}$ at $a=0.8$) because $|A^{i\eta}|\neq1$ there. The pole test therefore uses scaling with the distance to the pole, not absolute size.
\end{itemize}

\begin{thebibliography}{9}\small
\bibitem{BC} M.~M.~Bogdan, O.~V.~Charkina, arXiv:2502.14695, \emph{Low Temp. Phys.} 51 (2025).
\end{thebibliography}
\end{document}
